% ============================================================================
%  dag-generator — named TikZ styles (the canonical visual grammar).
%  Adapted from the SSVF/HUD-VASH methods supplement
%  (paper1-ssvf-hudvash/tex/tikz-causal-styles.tex). \input once after
%  \usepackage{tikz}. Figures use ONLY these named styles so the whole document
%  restyles from one place — and so the round-trip parser has a fixed grammar.
% ============================================================================
\usetikzlibrary{positioning, arrows.meta, fit, backgrounds,
                decorations.pathreplacing, calc, shapes.geometric}
\tikzset{
  % --- Nodes ---------------------------------------------------------------
  rv/.style     = {inner sep=2pt, minimum size=6mm},                 % random var: bare label
  cond/.style   = {draw, rectangle, inner sep=3.5pt, minimum size=7mm}, % conditioning box
  unobs/.style  = {text=gray!60!black, inner sep=2pt, minimum size=6mm}, % unobserved
  regime/.style = {draw=orange, text=orange, diamond, aspect=1.5,
                   inner sep=1pt, minimum size=6mm},                 % SWIG regime node
  fixed/.style  = {text=orange, inner sep=3pt, minimum size=6mm},    % set value (intervened)
  % --- Edges ---------------------------------------------------------------
  causal/.style = {-Latex},                                          % ordinary causal arrow
  det/.style    = {-Latex, double, double distance=1pt},             % deterministic relation
  reads/.style  = {-Latex, orange, dashed},                          % regime reads a var
  cfedge/.style = {-Latex, orange},                                  % regime -> cf descendant
  biasing/.style= {-Latex, red!70!black},                            % highlighted biasing path
  blocked/.style= {-Latex, gray, dashed},                            % blocked path
  bidirected/.style = {Latex-Latex, dashed, gray!60!black},          % unobserved common cause
  splitline/.style  = {orange, thick},                               % the node-split bar
}
% Counterfactual superscript marker (orange s), matching Y_t^s notation.
\providecommand{\cfs}{{\color{orange}s}}
% Colors just the intervention-value symbols inside potential-outcome notation,
% e.g. $H_1(\cfarg{a}, \cfarg{v_0})$.
\providecommand{\cfarg}[1]{{\color{orange}#1}}
% Node-split bar for SWIGs: \splitbar{<node>}{<fixed>} draws a vertical orange
% bar centered between the node and its fixed-value node.
\providecommand{\splitbar}[2]{%
  \draw[splitline] ($(#1.east)!0.5!(#2.west)+(0,3.5mm)$) -- ($(#1.east)!0.5!(#2.west)+(0,-3.5mm)$);}
